\section{Standing assumptions for our model}
\label{sec:assumptions}

We collect, once, the precise regularity conditions under which every result
below holds. They are mild and satisfied by every model in the companion package
(exponential kernels; power-law kernels with positive offset; non-negative,
locally integrable baselines). Each later result states any \emph{extra}
hypothesis explicitly.

\begin{standing}
\item \textbf{Simplicity and non-explosion.} $N$ is a simple temporal point
process on $[0,\infty)$: no two events coincide ($\Prob$-a.s.), and $N(t)<\infty$
a.s.\ for every finite $t$.
\item \textbf{Stochastic intensity.} $N$ admits an $(\Hist_t)$-predictable
conditional intensity $\lambda\ge 0$ with respect to its internal history
$\Hist_t=\sigma\{N(s):s\le t\}$; equivalently the compensator
$\Lambda(t)=\int_0^t\lambda$ is absolutely continuous, with $\int_0^T\lambda<\infty$
a.s.
\item \textbf{Kernel.} $\phi:\R\to[0,\infty)$ is measurable, causal ($\phi(u)=0$
for $u<0$), and locally bounded, with branching ratio
$n^{*}:=\int_0^\infty\phi=\kappa<\infty$. \emph{Subcritical} means $\kappa<1$.
\item \textbf{Baseline.} $s:[0,\infty)\to[0,\infty)$ is measurable and locally
integrable. Where a Dirac-comb (multi-impulse) baseline appears it is a
non-negative measure and convolutions with it are read distributionally.
\item \textbf{Finite mean intensity.} Wherever $\E[\lambda]$ is used,
$\sup_{t\le T}\E[\lambda(t)]<\infty$ --- automatic under (A3)-subcritical with
locally integrable $s$.
\item \textbf{Smooth, identifiable parameterisation.} For estimation, $\theta$
ranges over an open $\Theta\subseteq\R^p$; the map $\theta\mapsto\lambda(t;\theta)$
is twice continuously differentiable with $\lambda(t;\theta)>0$ on the data, and
the model is identifiable on $\Theta$.
\item \textbf{Multivariate.} In dimension $M$: $N=(N_1,\dots,N_M)$, intensity
$\lambda\in\R_+^M$, kernel \emph{matrix} $\Phi=(\phi_{mj})$, and \emph{branching
matrix} $G=(\int\phi_{mj})$. \emph{Stable} means the spectral radius
$\rho(G)<1$.
\end{standing}

\begin{intuition}
(A1)--(A2) say the process is well-behaved enough to have an instantaneous rate;
(A3)--(A4) say the two ingredients $\phi$ and $s$ are honest non-negative
functions with finite total excitation; (A5) keeps means finite; (A6) makes
maximum likelihood meaningful; (A7) lifts everything to $M$ interacting
components, where the single number $\kappa<1$ is replaced by the matrix condition
$\rho(G)<1$. Whenever a later theorem needs more --- boundedness of $\phi$,
subcriticality, existence of a transform --- it is flagged at that result.
\end{intuition}

\begin{remark}[Two parameterisations of the exponential kernel]
The package's event-time code writes $\phi(t)=\alpha e^{-\beta t}$, while the MBPP
theory uses $\phi(t)=\kappa\theta e^{-\theta t}$. They coincide with
$\alpha=\kappa\theta,\ \beta=\theta$, i.e.\ $\kappa=\alpha/\beta$ \emph{is} the
branching ratio. We use $(\kappa,\theta)$ throughout the theory.
\end{remark}
